\scanpage{263}
\mspage{44}
\margin{n° 227}
linéaires.

Soit $\unsure{\mathcal P|\mathcal S}$ un pencil. Un point $P'$ de $V$ est un \uline{point de base de} $\unsure{\mathcal P|\mathcal S}$ si $P'\in|X|$ pour tout $X\in\unsure{\mathcal P|\mathcal S}$.

THÉORÈME 3. -- \uline{Soit} $\unsure{\mathcal P|\mathcal S}$ \uline{un pencil de genre} $\geq1$. \uline{Alors} $\unsure{\mathcal P|\mathcal S}$ \uline{n'a pas de points de base qui soit simple sur} $V$.

\uline{Dém}. On va se servir d'un peu de variétés abéliennes ; Soit $k$ un corps de rationalité pour $V$, $\unsure{\mathcal P|\mathcal S}$ et $P'$, et soit $f:V\longrightarrow C$ la fonction définissant $\unsure{\mathcal P|\mathcal S}$. Soit $Q$ un point générique de $C/k$ et $P$ un point générique de $V$ tel que $f(P)=Q$. Alors on voit que $P'$ est une spécialisation de $P$ sur $k(Q)$, parce que $P'$ est un point dans le cycle premier rationnel ayant $P$ comme point générique sur $k(Q)$. Mais $P'$ étant rationnel sur $k$, on voit que la fonction $f$ ne peut être définie en $P'$. On peut supposer que $C$ est plongée dans sa Jacobienne. Ceci contredit le théorème 6 de VA.

\mspage{262}
